\section{Discussion}
\label{sec:discussion}

\paragraph{The central result is a dissociation.}
The central result is a dissociation: $\beta_k$ tracks whether the partner model is working, not whether the partner is worth engaging. A reliably hostile partner and a reliably cooperative one can both earn high $\beta_k$ --- what they share is predictability, not value. The tracked-only lesion confirms the mechanism: remove the link from $\beta_k$ to $\gamma_k$ and the behavioral gain disappears, even though partner-belief accuracy is preserved or slightly improved. The effect is upstream of inference quality, at the stage where existing beliefs are converted into action.

This supports the framing of affect as an augmentation layer over social inference rather than as a reward signal or a belief-accuracy mechanism. An agent with good ToM-like partner models can still be systematically under- or over-confident about acting on them; the affective precision channel calibrates this confidence independently of model content. The tracked-only lesion makes this experimentally concrete: it shows that ``better use of social beliefs'' and ``more accurate social beliefs'' are separable outcomes, and that the affective channel contributes to the former.

\paragraph{Metacognitive, not mentalistic.}
The mechanism is worth distinguishing from theory-of-mind inference. ToM approaches model the agent as maintaining representations of what partners believe, know, or intend---a second-order model of another agent's mental states \cite{FrithFrith2012,Pitliya2025ToMActiveInference}. The affective precision channel does not do this. It tracks confidence in the focal agent's own model of a partner's \emph{behavior}, not in representations of that partner's mental states. This makes it metacognitive rather than mentalistic: the signal is about how well the agent's own inference is performing, not about what the partner is thinking.

The two can complement each other. An agent with sophisticated ToM inference could still be poorly calibrated about when to act confidently on those inferences---particularly if the partner is reliable in some contexts and unreliable in others. Partner-local affective precision addresses exactly this: it allows action confidence to track the recent track record of each partner model, independently of the model's representational content. A natural extension is to combine ToM-style belief inference \cite{RuizSerra2025Factorised} with partner-local affective precision, so that the agent both models partner intentions and maintains a running estimate of how much those models currently justify decisive action.

\paragraph{The betrayal result reveals temporal structure.}
The abrupt-betrayal result is the strongest evidence that the mechanism is doing something real. If the affective precision channel were behaviorally inert, all betrayal variants would look the same. Instead, precision modulation lowers entropy, changes reallocation, and makes the agent more decisive---but around a stale model. The Hesp et al.\ architecture has two timescales: affective charge $\phi$ updates quickly from single-interaction prediction error, while the higher-order state $\beta$ integrates across many rounds. In stable environments this is fine --- $\beta$ accumulates evidence and converges. Under abrupt change there is a window in which $\phi_k$ has already registered surprise but $\beta_k$ still reflects pre-switch confidence. Action sharpens around a model that is no longer valid. This is not a bug in the mechanism so much as an exposure of its temporal assumption: it works when social regularities persist long enough for confidence and accuracy to stay aligned.

The gradual betrayal result confirms this interpretation directly: when behavioral change is slow enough that partner beliefs can track the shift before large confidence surplus accumulates, default precision modulation is neutral. The problem is not how much the affective channel responds, but how fast the social world changes relative to how fast beliefs reorganize.

The implication is that an effective affective precision channel in volatile environments would need a way to modulate the weight it places on recent charge signals based on how much the partner's behavior appears to be changing. This is structurally similar to hierarchical volatility tracking in belief-updating models \cite{Mathys2011HGF,Behrens2008SocialLearning}, but applied to action confidence rather than to learning rates. Whether these two levels of volatility sensitivity---for beliefs and for action precision---operate independently or are coupled in human social cognition is an open empirical question that this work motivates.

\paragraph{Limitations.}
Three design features bound the current claims. First, $\beta_k$ is an external auxiliary tracker; in the full Hesp et al.\ hierarchy it is a proper inferred hidden state, allowing the agent to have prior expectations about its own future affective states and to plan around anticipated changes in partner reliability. Making $\beta_k$ a variational variable would enable richer prospective regulation of action confidence. Second, partners are scripted rather than themselves being active-inference agents; the mechanism is tested in isolation from the reciprocal dynamics that characterize real social interaction. Third, the decisive test is a global-$\beta$ ablation: a condition that shares one affective precision estimate across all partners while preserving the same $\beta \to \gamma$ mapping. Without this, it remains possible that any precision modulation, not specifically partner-local precision, would produce similar benefits in the stable regimes. This comparison should be the highest priority for future work.

\paragraph{Future directions.}
The global-$\beta$ ablation is the most important near-term experiment. The predicted outcome is that a shared precision estimate would either over-suppress action confidence toward reliably predictable partners or under-suppress it toward uncertain ones, depending on the mix---but this should be tested. A second priority is a controlled shock-gradient study that varies betrayal abruptness parametrically, targeting the timescale boundary at which the mechanism transitions from neutral to harmful and testing the timing-mismatch account directly.

More broadly, the framework invites extension toward richer social environments. The trust game has no communication, reputation signaling, or coalition structure; all of these would introduce new sources of reliability evidence that the affective channel could in principle track. Testing with human participants would allow partner-local $\beta$ dynamics to serve as a computational model of individual differences in social confidence calibration---with faster, higher-gain channels predicted to show stronger post-cue reallocation but greater sensitivity to abrupt social change. Finally, combining this mechanism with ToM-style belief inference \cite{Pitliya2025ToMActiveInference} would allow future work to study whether metacognitive action-confidence signals and mentalistic partner-state models interact during social decision-making, and how they trade off under resource or time constraints.

\section{Conclusion}

Predicting what a partner will do and deciding how confidently to act on that prediction are separable problems. This paper shows that affect, formalized as a partner-specific estimate of social model fitness, can address the second problem. In repeated trust-game simulations, partner-local affective precision tracks partner predictability rather than partner value, improves performance through action sharpening in decision-open regimes, and reorganizes partner engagement around reliability before aggregate payoff separates. When partners change abruptly, the same mechanism sharpens action around a model that has not yet caught up---a failure that reveals the timing structure of the mechanism rather than a flaw in its steady-state design. The value of calibrating action confidence to social model fitness is real and the limits are principled: they tell us something about the timescales on which adaptive social behavior depends.
